\documentclass[11pt]{article}
\usepackage[a4paper,margin=24mm]{geometry}
\usepackage{amsmath,amssymb,amsthm,mathtools}
\usepackage{booktabs,tabularx,array}
\usepackage{graphicx}
\usepackage{placeins}
\usepackage[hidelinks]{hyperref}
\usepackage{microtype}

\graphicspath{{docs/paper/figures/}}
\newcommand{\FEMPS}{\textnormal{\textsc{femps}}}
\newcommand{\claimid}[1]{\nolinkurl{#1}}
\DeclareMathOperator{\Tr}{Tr}
\newcolumntype{Y}{>{\raggedright\arraybackslash}X}

\title{Restricted Nonbranching Functional Exterior Matrix Product States\\
for Continuous Interacting Fermions}
\author{Author information withheld for review}
\date{}

\begin{document}
\maketitle

\begin{abstract}
Functional tensor networks expand each continuous particle coordinate in an
orthonormal basis, contract the resulting operator matrices, and optimize a
finite coefficient representation by automatic differentiation.  Extending
this construction to identical fermions is difficult: an ordinary
particle-site tensor train carries exchange-statistics ranks, whereas generic
compact exterior cores do not admit an assumed efficient exact contraction.
We therefore study a restricted functional exterior matrix product state
(\FEMPS) with one globally conserved, nonbranching virtual path.  The state is
exactly a sum of $K$ nonorthogonal Slater determinants in a continuous
functional basis, embedded in the matrix-wedge definition.  Norms, one-body
and factorized two-body energies, and reverse-mode gradients are evaluated
through batched $K^2$ determinant transitions, with inverse-based fast formulas
and a singular-safe minor fallback.  A conditioned generalized eigenproblem removes
the linear determinant amplitudes during orbital optimization.

Numerical evidence is reported for harmonic and soft-Coulomb fermions.  At
$N=4$, independent $D$ and $K$ axes, three-seed stability, and a
$D=8\to12$ continuation pass direct finite-basis configuration-interaction
(CI), variance, norm, antisymmetry, time, and memory checks.  At
$N=6,D=10,K=4$, three blind runs have CI errors
$2.49\times10^{-4}$--$4.03\times10^{-4}$, maximum variance
$3.09\times10^{-3}$, and zero structural antisymmetry residual.  Independent
$N=6$ axes further give monotone energies for $K=1,2,4$ at $D=10$ and
$D=8,10,12$ at $K=4$; the final $D=12$ CI error is
$1.05\times10^{-4}$.  A preregistered truth-free $D=12$ adaptive continuation
to $K=6$ further reduces that error to $3.20\times10^{-5}$ and the variance to
$3.73\times10^{-4}$, while beating a same-budget cold $K=6$ control.  At fixed $D=10,K=4$ the ordinary particle-TT center rank is 24 for $N=4$ and 80 for
$N=6$, while the \FEMPS{} correlation control remains $K=4$; matched kernel
measurements expose the corresponding polynomial cost.  The method is a
restricted nonorthogonal multideterminant algorithm, not an efficient generic
\FEMPS{} contraction.  Direct CI remains faster in the present small truth
spaces, and no $N=8$, asymptotic-scaling, DMRG-superiority, or categorical
novelty claim is made.
\end{abstract}

\section{Introduction}

Hong et al. introduced a functional tensor-network route to continuous
many-body Schr\"odinger equations \cite{Hong2022FunctionalTN}.  For an
orthonormal basis $\{\phi_p\}_{p=1}^{D}$, one writes
\begin{equation}
 \Psi(x_1,\ldots,x_N)=
 \sum_{p_1,\ldots,p_N=1}^{D} C_{p_1\cdots p_N}
 \prod_{i=1}^{N}\phi_{p_i}(x_i).
 \label{eq:functional-expansion}
\end{equation}
Orthonormality converts continuum norms and local differential or
multiplicative operators into finite tensor contractions.  Hong et al.
represented $C$ by an MPS and optimized the Rayleigh quotient by automatic
differentiation, leaving fermionic tensor networks as an open extension.

For identical fermions, $C$ lies in $\Lambda^N V_D$ rather than a generic
$V_D^{\otimes N}$.  An ordinary particle-site TT can represent this alternating
tensor, but its exact bonds include exchange multiplicities even for a single
Slater determinant.  Occupation-number MPS/DMRG, Hilbert-space MPS
\cite{ChanSharma2011DMRG,LiChan2016HSMPS}, ordered first-quantized MPS
\cite{LiWaintal2026FirstQuantization}, and an exterior particle-coordinate
network are distinct factorizations and must retain their own names.

Exterior algebra enforces antisymmetry structurally, but compact exterior
inputs do not by themselves imply tractable exact contraction.  A companion
structural analysis conditionally obstructs generic exact squared-norm
contraction already at fixed small bond.  Here that result is treated as an
algorithm-design constraint: we impose a nonbranching global virtual label and
test the resulting restricted class as a continuous variational solver.

The contribution is deliberately bounded:
\begin{enumerate}
 \item an exact matrix-wedge embedding of a $K$-term nonorthogonal Slater sum;
 \item polynomial value and reverse-mode determinant-transition contractions,
       including singular-safe validation;
 \item a conditioned variable-projection optimizer with versioned checkpoints,
       residuals, and reproducibility records;
 \item controlled $N=2,4,6$ harmonic and soft-Coulomb benchmarks; and
 \item a measured separation between the exterior exchange carrier and the
       correlation multiplicity $K$, together with its cost.
\end{enumerate}
Finite sums of nonorthogonal Slaters and geminal states have substantial prior
art \cite{BeylkinMohlenkampPerez2008SlaterSums,UemuraKasamatsuSugino2015AGPCI}.
Accordingly, we do not claim a new determinant state class.  The method-level
content is the restricted exterior/functional integration, its contraction and
AD implementation, and its controlled continuum-basis audit.

\section{Functional exterior representation}

\subsection{Finite functional basis}

Let $V_D=\operatorname{span}\{\phi_1,\ldots,\phi_D\}$ with
$\langle\phi_p,\phi_q\rangle=\delta_{pq}$.  A normalized decomposable state is
\begin{equation}
 \Phi[U]=u_1\wedge\cdots\wedge u_N,\qquad
 U=(u_1,\ldots,u_N)\in\mathbb C^{D\times N},\quad U^\dagger U=I_N.
 \label{eq:slater}
\end{equation}
The normalized alternating embedding into labelled particle coordinates is
\begin{equation}
 \mathcal J_N(u_1\wedge\cdots\wedge u_N)
 =\frac{1}{\sqrt{N!}}\sum_{\pi\in S_N}\operatorname{sgn}(\pi)
 u_{\pi(1)}\otimes\cdots\otimes u_{\pi(N)}.
 \label{eq:alternating-embedding}
\end{equation}
Antisymmetry is therefore exact before any numerical contraction.

\subsection{Nonbranching diagonal-path \FEMPS}

A matrix-wedge \FEMPS{} propagates virtual matrix indices while multiplying
physical entries by exterior product.  We restrict every core to be diagonal
in one global path label $a=1,\ldots,K$.  With boundary amplitudes $c_a$, the
state becomes
\begin{equation}
 \Psi_{K,D,N}=\sum_{a=1}^{K}c_a\,
 u_{a1}\wedge\cdots\wedge u_{aN}
 =\sum_{a=1}^{K}c_a\Phi_a.
 \label{eq:diagonal-path}
\end{equation}
No core can change $a$, hence there are $K$ physical paths rather than
$K^{N-1}$ branching virtual histories.  Equation~\eqref{eq:diagonal-path} is an
exact subclass of the original matrix-wedge definition and remains first
quantized in the continuous coordinates of Eq.~\eqref{eq:functional-expansion}.

$K=1$ is a single Slater.  Increasing $K$ enlarges a nonorthogonal secant span
and can recover correlation.  We call $K$ the \emph{correlation multiplicity}
of this restricted parametrization; it is not a canonical entanglement
spectrum.  The stored nonlinear state has $KDN$ complex orbital scalars before
gauge removal.

\section{Exact restricted contractions}

Let $U_a\in\mathbb C^{D\times N}$ collect the orbitals of $\Phi_a$ and define
the transition overlap
\begin{equation}
 S_{ab}=U_a^\dagger U_b,\qquad
 \langle\Phi_a|\Phi_b\rangle=\det S_{ab}.
 \label{eq:overlap}
\end{equation}
All norm and Hamiltonian matrices therefore have only $K^2$ transition pairs.

\subsection{One-body transitions}

For $\hat H_1=\sum_i h(i)$, let $X_{ab}=U_a^\dagger hU_b$.  If $S_{ab}$ is
well conditioned,
\begin{equation}
 \langle\Phi_a|\hat H_1|\Phi_b\rangle
 =\det(S_{ab})\Tr(S_{ab}^{-1}X_{ab}).
 \label{eq:one-body-transition}
\end{equation}
Near singularity, the implementation replaces one column of $S_{ab}$ at a
time and sums the resulting determinants.  This cofactor formula is valid even
when $S_{ab}$ is singular and is also an independent value/gradient reference.

\subsection{Factorized two-body transitions}

The admitted pair operator is written
\begin{equation}
 \hat V=\frac12\sum_{\ell=1}^{L}w_\ell\sum_{i<j}
 \left[L_\ell(i)R_\ell(j)+R_\ell(i)L_\ell(j)\right].
 \label{eq:factorized-pair}
\end{equation}
For projected insertions
$A_{ab}^{\ell}=U_a^\dagger L_\ell U_b$ and
$B_{ab}^{\ell}=U_a^\dagger R_\ell U_b$, the well-conditioned mixed transition is
\begin{equation}
 \det(S_{ab})\left[
 \Tr(S_{ab}^{-1}A_{ab}^{\ell})\Tr(S_{ab}^{-1}B_{ab}^{\ell})
 -\Tr(S_{ab}^{-1}A_{ab}^{\ell}S_{ab}^{-1}B_{ab}^{\ell})
 \right].
 \label{eq:two-body-transition}
\end{equation}
The singular-safe route replaces two distinct columns and evaluates
$K^2LN(N-1)$ determinants.  The auto path selects inverse formulas only after a
conditioning gate.  Well-conditioned $K^2$ pairs and all $L$ factors are
evaluated as batched solves; only singular pairs enter the column-replacement
loop.  Both branches remain differentiable under PyTorch AD, while the former
pairwise implementation is retained as an independent reference.

\subsection{Energy and complexity}

With overlap and Hamiltonian matrices $S,H\in\mathbb C^{K\times K}$,
\begin{equation}
 E(c,U)=\frac{c^\dagger H(U)c}{c^\dagger S(U)c}.
 \label{eq:rayleigh}
\end{equation}
The stored orbitals require $O(KDN)$ memory.  Dense one-body and factor actions
cost $O(KD^2N)$ and $O(KLD^2N)$, followed by $K^2$ transition algebra.  A
conservative well-conditioned bound is polynomial in
$K^2L(D^2N+N^3)$.  The validation fallback performs
$K^2LN(N-1)$ size-$N$ determinants and is slower but still polynomial.  Neither
route enumerates all virtual paths or materializes $D^N$ coefficients.

\section{Variable-projection solver and scientific contract}

Each $U_a$ is kept in a Stiefel gauge by reduced QR, with the diagonal phase of
$R$ fixed deterministically.  At every nonlinear orbital step, amplitudes are
eliminated by the conditioned generalized Hermitian problem
\begin{equation}
 H(U)c=E\,S(U)c.
 \label{eq:generalized-eigenproblem}
\end{equation}
Overlap directions below a relative threshold are discarded and the retained
rank, raw/retained condition number, and generalized residual are reported.
Adam with a cosine learning-rate schedule is followed by optional L-BFGS
refinement.  Reverse-mode derivatives propagate through QR, transitions, and
the generalized eigenvalue.

The public result schema is version 2 and the checkpoint schema is version 1.
Resume rejects a changed configuration, operator identity, or schema.  Every
call reports the structural antisymmetry residual, norm error, energy variance
when exterior truth is enabled, elapsed time, sampled process RSS, conditioning,
and exact operation counts.  If $D^N$ lies below a declared cap, a validation-
only particle tensor is built and its adjacent-transposition residual is also
reported.  A null materialized residual means that the registered exponential
validator was disabled; it never replaces the structural residual.

The machine-readable reproduction manifest records SHA-256 hashes, commands,
seeds, tolerances, evidence labels, and independent verifiers.  Claim
identifiers used below refer to that manifest.

\section{Models and reference calculations}

We use harmonic-oscillator functions as the continuous basis.  The first model
is an analytically separable harmonically interacting fermion system, used for
$N=2$ validation and the initial $N=4$ ladder.  The nonquadratic model is
\begin{equation}
 \hat H=\sum_{i=1}^{N}\left(-\frac12\partial_{x_i}^2+\frac12x_i^2\right)
 +\sum_{i<j}\frac{1}{\sqrt{(x_i-x_j)^2+1}}.
 \label{eq:soft-coulomb}
\end{equation}
Soft-Coulomb integrals are projected by $Q=128$ Gauss--Hermite quadrature.  The
production operator uses a physical $(p,r)\mid(q,s)$ SVD of the dense pair
tensor; $D=10$ and 12 retain ranks 19 and 23 with relative reconstruction
errors $1.34\times10^{-15}$ and $3.85\times10^{-15}$.

Direct references use the unfactorized four-index quadrature tensor and the
Slater--Condon Hamiltonian in $\Lambda^N V_D$.  CI is validation-only and never
initializes a \FEMPS{} run.  At $N=6,D=10$ and 12 the exterior dimensions are
$\binom{10}{6}=210$ and $\binom{12}{6}=924$, so exact diagonalization remains stronger than an
approximate second-quantized DMRG control.  DMRG is therefore deferred rather
than relabeled as \FEMPS{}.

\section{Numerical evidence}

All floating-point results in this section are numerical evidence.  The full
manifest contains eleven independently verified artifacts.

\subsection{\texorpdfstring{$N=2$}{N=2}: analytic and finite-basis checks}

The noninteracting $N=2,D=6,K=1$ state gives $E=2$ and zero variance exactly
within complex128 arithmetic.  For harmonic coupling $\kappa=0.35$, increasing
$K=1,2,4$ at $D=6$ lowers the energy, while $D=4,6,8$ at $K=4$ decreases the
continuum-reference error.  The final $D=8,K=4$ energy is
$2.4557706074300585$, with continuum-reference error
$9.886\times10^{-6}$, same-basis error $9.797\times10^{-6}$, variance
$8.423\times10^{-5}$, and zero antisymmetry residual
(claim \claimid{n2_harmonic_ladder}).

The same registered ladder embeds the noninteracting $N=4$ Slater at
$D=6,K=1$, yielding $E=8$ with zero variance and zero materialized
antisymmetry residual.  For the interacting harmonic model at $\kappa=0.35$,
three blind $N=4,D=6,K=4$ runs agree within $1.782\times10^{-12}$ and have
maximum same-basis error $1.835\times10^{-12}$ and maximum variance
$2.272\times10^{-11}$.  The $K=1,2,4$ energy is nonincreasing at $D=6$ and
the continuum-reference error decreases over $D=5,6,7$ at $K=4$ (claim
\claimid{n4_harmonic_e4_closure}).

\subsection{\texorpdfstring{$N=4$}{N=4}: nonquadratic $D$ and $K$ control}

For Eq.~\eqref{eq:soft-coulomb}, Table~\ref{tab:n4-k} shows the independent
$D=6$ correlation axis.  Increasing $K$ reduces the same-basis CI error by
more than four orders of magnitude.

\begin{table}[t]
\centering
\caption{\texorpdfstring{$N=4,D=6$}{N=4,D=6} soft-Coulomb correlation axis.}
\label{tab:n4-k}
\begin{tabular}{rrrr}
\toprule
$K$ & Energy & Error versus CI & Variance \\
\midrule
1 & 11.025279022118180 & $1.441\times10^{-3}$ & $6.360\times10^{-3}$ \\
2 & 11.023866678350620 & $2.897\times10^{-5}$ & $2.009\times10^{-4}$ \\
4 & 11.023837776253766 & $6.305\times10^{-8}$ & $3.147\times10^{-7}$ \\
\bottomrule
\end{tabular}
\par\smallskip{\scriptsize Evidence:
\claimid{n4_soft_coulomb_transferability}.}
\end{table}

At fixed $K=4$, the absolute energy difference from a pre-existing $D=14$
finite-basis numerical reference decreases from $7.549\times10^{-4}$ at $D=6$
to $9.494\times10^{-5}$ at $D=12$.  This is not a continuum bound.  Direct
same-basis CI errors are $1.042\times10^{-4}$ and $8.314\times10^{-5}$ at
$D=10,12$, with variances below $10^{-3}$ and zero structural/materialized
antisymmetry residuals (claim
\claimid{n4_soft_coulomb_basis_extension}).

Blind growth from $D=12,K=4$ to $K=5$ lowers the same-basis CI error by 44.35\%
to $4.627\times10^{-5}$ and the variance to $5.291\times10^{-4}$, while the
measured time grows from 92.77 to 141.91 seconds (claim
\claimid{n4_soft_coulomb_high_basis_correlation}).  Three blind $D=6$ starts
and three truth-free $D=6\to8$ continuations independently pass the registered
stability gates.

\begin{figure}[t]
 \centering
 \includegraphics[width=\textwidth]{femps-convergence-summary.pdf}
 \caption{Restricted \FEMPS{} convergence from manifest-hashed artifacts.
 Panel (a) is the N4 same-basis correlation error, panel (b) is the N4
 functional-basis trend relative to a finite $D=14$ numerical reference, and
 panel (c) separates N6 K1, K1-to-K4 continuation, and blind K4 starts.  All
 values are numerical evidence.}
 \label{fig:convergence}
 \par\smallskip{\scriptsize Evidence:
 \claimid{n4_soft_coulomb_transferability},
 \claimid{n4_soft_coulomb_basis_extension},
 \claimid{n4_soft_coulomb_high_basis_correlation},
 \claimid{n6_soft_coulomb_pilot}, and
 \claimid{n6_soft_coulomb_multiseed}.}
\end{figure}

\subsection{\texorpdfstring{$N=6$}{N=6}: feasibility and blind stability}

The resource-capped $N=6,D=10$ pilot first optimizes a blind $K=1$ Slater and
then preserves it while adding three seeded blind determinants.  The CI energy
is $25.049639839832263$.  The $K=1$ error is $2.603\times10^{-3}$; $K=4$
reduces it to $1.382\times10^{-4}$ with variance $1.197\times10^{-3}$, zero
materialized antisymmetry residual, 85.59-second runtime, and 798 MB sampled
peak RSS (claim \claimid{n6_soft_coulomb_pilot}).

Three additional, genuinely blind K4 starts use identical 160-step Adam and
80-step L-BFGS budgets.  Table~\ref{tab:n6-seeds} reports all runs rather than
selecting the best initialization.

\begin{table}[t]
\centering
\caption{Blind \texorpdfstring{$N=6,D=10,K=4$}{N=6,D=10,K=4} stability.}
\label{tab:n6-seeds}
\begin{tabular}{rrrrr}
\toprule
Seed & CI error & Variance & Condition number & Time (s) \\
\midrule
31 & $2.494\times10^{-4}$ & $2.169\times10^{-3}$ & 3.394 & 112.14 \\
37 & $2.932\times10^{-4}$ & $2.334\times10^{-3}$ & 3.209 & 105.06 \\
43 & $4.031\times10^{-4}$ & $3.094\times10^{-3}$ & 4.424 & 105.92 \\
\bottomrule
\end{tabular}
\par\smallskip{\scriptsize Evidence: \claimid{n6_soft_coulomb_multiseed}.}
\end{table}

The energy spread is $1.537\times10^{-4}$; all four overlap directions are
retained, norm and structural antisymmetry residuals are zero, and every point
stays below 883 MB sampled peak RSS.  Seed 31 additionally materializes the
one-million-coefficient tensor and gives zero antisymmetry residual.

Independent functional-basis and correlation axes are reported in
Table~\ref{tab:n6-dk}.  Every point uses 160 Adam and 80 L-BFGS steps.  The
$D=12,K=4$ point takes 107.66 seconds and 974,974,976 sampled RSS bytes.  Its
structural antisymmetry residual is zero; no $12^6$ particle tensor is built.

\begin{table}[t]
\centering
\caption{Independent \texorpdfstring{$N=6$}{N=6} soft-Coulomb convergence axes.}
\label{tab:n6-dk}
\begin{tabular}{llrrr}
\toprule
Axis & Point & Energy & Error versus CI & Variance \\
\midrule
$K$ at $D=10$ & $K=1$ & 25.052242782725 & $2.603\times10^{-3}$ & $1.478\times10^{-2}$ \\
 & $K=2$ & 25.050276338618 & $6.365\times10^{-4}$ & $4.695\times10^{-3}$ \\
 & $K=4$ & 25.049825287522 & $1.854\times10^{-4}$ & $1.749\times10^{-3}$ \\
$D$ at $K=4$ & $D=8$ & 25.051264850894 & $3.196\times10^{-5}$ & $3.016\times10^{-4}$ \\
 & $D=10$ & 25.049825287522 & $1.854\times10^{-4}$ & $1.749\times10^{-3}$ \\
 & $D=12$ & 25.049471144618 & $1.047\times10^{-4}$ & $1.122\times10^{-3}$ \\
\bottomrule
\end{tabular}
\par\smallskip{\scriptsize Evidence: \claimid{n6_independent_dk_convergence}.}
\end{table}

Both the \FEMPS{} and direct-CI absolute energies decrease over $D=8,10,12$.
The error relative to each same-basis CI result is not monotone because fixed
$K=4$ must approximate new correlation directions as the basis grows.  This
is an interpretable $D/K$ tradeoff, not a continuum bound.  At $D=12,K=4$ the
state stores 292 complex \FEMPS{} scalars; its ordinary particle-TT ranks are
$(12,60,80,60,12)$ and require 132,768 scalars (claim
\claimid{n6_independent_dk_convergence}).

ADR 0023 next freezes a truth-free adaptive $D=12$ continuation.  At each
growth step, exactly 32 seeded Slater candidates are ranked only by the
factorized determinant-transition generalized eigenvalue and balanced overlap
conditioning.  Dense CI is constructed only after the adaptive K5, adaptive
K6, and cold K6 optimizations are frozen.  Table~\ref{tab:n6-adaptive} reports
the result.

\begin{table}[t]
\centering
\caption{Truth-free adaptive \texorpdfstring{$N=6,D=12$}{N=6,D=12} growth.}
\label{tab:n6-adaptive}
\begin{tabular}{lrrrr}
\toprule
State & Energy & Error versus CI & Variance & Condition \\
\midrule
Source $K=4$ & 25.049471144618 & $1.047\times10^{-4}$ & $1.122\times10^{-3}$ & 2.111 \\
Adaptive $K=5$ & 25.049434533568 & $6.812\times10^{-5}$ & $7.747\times10^{-4}$ & 2.440 \\
Adaptive $K=6$ & 25.049398437461 & $3.202\times10^{-5}$ & $3.726\times10^{-4}$ & 2.760 \\
Cold $K=6$ & 25.049642546632 & $2.761\times10^{-4}$ & $2.659\times10^{-3}$ & 8.039 \\
\bottomrule
\end{tabular}
\par\smallskip{\scriptsize Evidence: \claimid{n6_adaptive_correlation_growth}.}
\end{table}

Adaptive K6 lowers the source energy by $7.271\times10^{-5}$ and lies
$2.441\times10^{-4}$ below the same-budget cold K6.  Its ordinary particle-TT
ranks are $(12,66,120,66,12)$, requiring 209,376 scalars, while the restricted
\FEMPS{} state stores 438 orbital and amplitude scalars.  This is a structural
representation comparison, not a runtime advantage.  K5, adaptive K6, and
cold K6 take 6.72, 6.24, and 6.28 CPU seconds after vectorization; the 107.66 s
source record predates Phase 33 and is not a matched timing comparator.  The
K6 pruning gate does not trigger (claim
\claimid{n6_adaptive_correlation_growth}).

\section{Exchange carrier, correlation multiplicity, and cost}

At fixed $D=10,K=4,L=19$, the exact structural counts grow as shown in
Table~\ref{tab:cost-counts}.  The $N=6/N=4$ ratio is 1.5 for stored orbitals and
one-body determinants and 2.5 for two-body minors.

\begin{table}[t]
\centering
\caption{Matched production structure.}
\label{tab:cost-counts}
\begin{tabular}{lrrr}
\toprule
Quantity & $N=4$ & $N=6$ & Ratio \\
\midrule
Stored orbital scalars & 160 & 240 & 1.5 \\
Transition pairs & 16 & 16 & 1.0 \\
One-body determinants & 64 & 96 & 1.5 \\
Two-body determinants & 3,648 & 9,120 & 2.5 \\
Enumerated virtual paths & 0 & 0 & -- \\
\bottomrule
\end{tabular}
\par\smallskip{\scriptsize Evidence: \claimid{matched_n4_n6_transition_cost}.}
\end{table}

The singular-safe minor value/gradient timings grow by factors 2.115/2.232,
close to the structural count after fixed overhead.  The inverse-fast path
grows by 1.239/1.024 over these two small points and accelerates the N6 minor
path by 7.27/8.93.  These are two-point CPU measurements, not asymptotic fits.
Auto and minor transition matrices agree within $1.07\times10^{-14}$.

\FloatBarrier

Phase 33 batches transition pairs and factor axes while retaining that
pairwise implementation as a reference.  Table~\ref{tab:vectorized-backend}
uses one fixed $N=6,D=10,K=4,L=19$ state and one matched blind optimization.

\begin{table}[t]
\centering
\caption{Vectorized transition and matched backend audit.}
\label{tab:vectorized-backend}
\small
\begin{tabular}{lrrr}
\toprule
Path & FB kernel (s) & Reference speedup & 160/80 solve (s) \\
\midrule
CPU pairwise reference & 0.28038 & 1.0 & -- \\
CPU batched & 0.00803 & 34.9 & 5.08 \\
Blackwell batched & 0.01814 & 15.5 & 11.96 \\
\bottomrule
\end{tabular}
\par\smallskip{\scriptsize Evidence:
\claimid{n6_vectorized_transition_backend}.}
\end{table}

CPU batched/reference Hamiltonians differ by $8.90\times10^{-16}$ and orbital
gradients by $4.69\times10^{-13}$.  CPU/Blackwell Hamiltonians differ by
$1.43\times10^{-14}$ and gradients by $3.68\times10^{-12}$.  The matched
optimizations give identical energy $25.050223374042$; Blackwell finishes well
inside 600 seconds with 20.7 MB peak device allocation.  CPU remains faster by
a factor 2.35 for the full registered solve, so Blackwell is an admitted backend
but not an acceleration claim.  Seed 3304 misses the separate Phase 29 CI-error
quality control and is used only for backend parity; it does not replace the
Phase 32 physics evidence.

\begin{figure}[htbp]
 \centering
 \includegraphics[width=0.72\textwidth]{femps-structure-cost-summary.pdf}
 \caption{Matched N4/N6 kernel timings and ordinary particle-TT center ranks
 at fixed \FEMPS{} $K=4$.  The exterior representation keeps the correlation
 multiplicity fixed while the labelled-particle exchange carrier grows.  The
 plot does not imply runtime superiority over CI.}
 \label{fig:structure-cost}
 \par\smallskip{\scriptsize Evidence:
 \claimid{matched_n4_n6_transition_cost}.}
\end{figure}

\FloatBarrier

For $N=4$ the ordinary particle-TT ranks are $(10,24,10)$ for the K4 state and
$(10,45,10)$ for dense CI.  At $N=6$ they are
$(10,45,80,45,10)$ and $(10,45,120,45,10)$.  The \FEMPS{} nonlinear
correlation multiplicity remains $K=4$ even though the ordinary center rank
increases from 24 to 80.  Exterior algebra does not remove exchange support;
it prevents that support from being conflated with the adjustable correlation
count.

\section{Relation to other representations}

The diagonal-path state is a nonorthogonal selected-CI-like Slater expansion,
and determinant sums are established variational objects
\cite{BeylkinMohlenkampPerez2008SlaterSums}.  Finite AGP/AGP-CI methods likewise
use polynomial pair-state transitions \cite{UemuraKasamatsuSugino2015AGPCI}.
The present implementation differs in its explicit derivation as a restricted
matrix-wedge \FEMPS{}, direct use of the Hong et al. continuous
functional-operator layer, and mandatory $D/K$/symmetry/cost audit.  These
differences do not establish categorical novelty of the underlying state
family.

An ordinary particle TT \cite{Oseledets2011TT} is an exact comparator in the
labelled coordinate tensor product.  Its ranks include the flat exchange
spectrum of a Slater determinant, a standard fermionic reduced-density
structure \cite{Coleman1963FermionRDM}.  Occupation-space quantum-chemistry
DMRG \cite{ChanSharma2011DMRG} and Hilbert-space MPS
\cite{LiChan2016HSMPS} operate in a different factorization.  The ordered
first-quantized route of Li--Waintal \cite{LiWaintal2026FirstQuantization}
removes labelled-particle redundancy by restricting to an ordered sector; it
is an external first-quantized comparator, not a \FEMPS{} backend.

\section{Limitations and failure conditions}

The strongest limitation is scope.  Equation~\eqref{eq:diagonal-path} forbids
virtual branching and therefore does not solve generic matrix-wedge
contraction.  Expressivity is that of a $K$-Slater secant expansion; difficult
states may require rapidly growing $K$.  The measured cost is explicitly
quadratic in $K$ and linear in the two-body factor rank $L$ before orbital and
determinant costs.

Direct CI is faster in the current truth spaces, which extend to
$\binom{12}{6}=924$.  The evidence demonstrates correctness, optimization
stability, three-point N6 basis/correlation convergence, and a structural
tradeoff, not an end-to-end speed or memory advantage.  Adaptive K growth has
one registered candidate-pool lineage and one cold control; it does not yet
establish pool-seed stability or a universal stopping rule.  Three basis points and
two particle counts cannot establish asymptotic scaling; no $N=8$ diagonal-path
calculation is admitted.  DMRG has been deferred
because it would be weaker than exact CI in the present finite spaces; no
FEMPS--DMRG performance inference follows.

The implementation uses exact deterministic contraction for the restricted
class.  A future truncation or stochastic estimator for generic cores would
require a separate error, variance, failure-probability, and antisymmetry gate.
No floating-point experiment in this paper is a theorem.

\section{Conclusion}

A globally conserved nonbranching path turns a matrix-wedge functional
exterior network into a $K$-term nonorthogonal Slater sum.  This restriction is
strong enough to avoid known generic contraction obstructions while preserving
first quantization, continuous functional bases, exact antisymmetry, and an
adjustable multideterminant correlation span.  Determinant transitions provide
exact norms, one- and two-body energies, and AD gradients without virtual-path
or particle-tensor enumeration in production.

Controlled N2/N4/N6 calculations establish a functioning restricted solver:
energies improve systematically in $D$ and $K$ on registered axes, blind
interacting runs are stable, and every accepted approximation/error/resource
quantity is explicit.  Truth-free fixed-pool growth at N6,D12 converts the
correlation control into a measurable solver improvement, while its single-pool
scope remains explicit.  The ordinary particle-TT comparison makes the intended
separation visible: exchange support grows while $K$ remains the chosen
correlation multiplicity.  At the same time, current CI references are faster
and the method is close to established nonorthogonal determinant expansions.
The result is therefore a bounded algorithmic recovery of \FEMPS{}, not a
generic scalability claim.

\appendix
\section{Reproduction and evidence mapping}

The canonical short audit is
\begin{verbatim}
python scripts/build_phase30_reproduction_manifest.py
python scripts/verify_phase30_reproduction_manifest.py
python scripts/build_phase30_method_figures.py
python scripts/verify_phase30_method_figures.py
python scripts/verify_phase31_method_manuscript.py
python scripts/build_phase31_method_manuscript.py
python -m pytest -q
\end{verbatim}
Long benchmark commands, seeds, tolerances, scientific boundaries, artifact
hashes, and verifier modules are stored in
\texttt{docs/experiments/results/phase30\_reproduction\_manifest.json}.
Checkpoints are ignored regeneration artifacts and are not committed evidence.

\begin{table}[h]
\centering
\caption{Manuscript numerical claims and manifest identifiers.}
\begin{tabularx}{\textwidth}{lY}
\toprule
Identifier & Scope \\
\midrule
\claimid{n2_harmonic_ladder} & N2 E1/E2 and noninteracting N4 E3 \\
\claimid{n4_harmonic_e4_closure} & Interacting harmonic N4 closure \\
\claimid{n4_soft_coulomb_transferability} & N4 nonquadratic multiseed and D/K axes \\
\claimid{n4_soft_coulomb_basis_extension} & N4 D8-to-D12 continuation \\
\claimid{n4_soft_coulomb_high_basis_correlation} & N4 D12 K4-to-K5 growth \\
\claimid{n6_soft_coulomb_pilot} & N6 K1-to-K4 resource-capped pilot \\
\claimid{n6_soft_coulomb_multiseed} & N6 three-blind-seed stability \\
\claimid{matched_n4_n6_transition_cost} & Fixed-D/K/L N4-to-N6 cost audit \\
\claimid{n6_independent_dk_convergence} & N6 independent D8/D10/D12 and K1/K2/K4 axes \\
\claimid{n6_vectorized_transition_backend} & Batched CPU/Blackwell parity and matched solve \\
\claimid{n6_adaptive_correlation_growth} & Truth-free N6,D12 K4-to-K6 growth and cold K6 control \\
\bottomrule
\end{tabularx}
\end{table}

\begingroup
\footnotesize
\bibliographystyle{unsrt}
\bibliography{references/references}
\endgroup

\end{document}
